% ============================================================================
%  Score-Following: Алгоритмы отслеживания партитуры в реальном времени
%  MusicTech --- Центральный Университет
%  Автор: Никита Новицкий, 2026
% ============================================================================
\documentclass[a4paper,11pt]{article}

% --- Кодировка и язык ---
\usepackage[utf8]{inputenc}
\usepackage[T2A]{fontenc}
\usepackage[russian]{babel}

% --- Геометрия ---
\usepackage[left=1.8cm,right=1.8cm,top=2.2cm,bottom=2.2cm]{geometry}

% --- Математика ---
\usepackage{amsmath,amssymb,amsfonts}

% --- Графика и цвет ---
\usepackage{graphicx}
\usepackage[dvipsnames,svgnames]{xcolor}
\usepackage{tikz}
\usepackage{pgfplots}
\pgfplotsset{compat=1.18}

% --- TikZ библиотеки ---
\usetikzlibrary{
  arrows.meta,
  positioning,
  shapes.geometric,
  calc,
  fit,
  backgrounds,
  decorations.pathreplacing
}

% --- Таблицы ---
\usepackage{booktabs}
\usepackage{tabularx}

% --- tcolorbox ---
\usepackage[most]{tcolorbox}

% --- Гиперссылки ---
\usepackage[
  colorlinks=true,
  linkcolor=accentblue,
  urlcolor=accentblue!80!black,
  citecolor=accentblue!70!black
]{hyperref}

% --- Колонтитулы ---
\usepackage{fancyhdr}

% --- Списки ---
\usepackage{enumitem}

% --- Плавающие объекты ---
\usepackage{float}

% --- Заголовки секций ---
\usepackage{titlesec}

% --- Интервалы ---
\usepackage{setspace}

% --- Листинги кода ---
\usepackage{listings}

% --- Подписи ---
\usepackage{caption}

% ============================================================================
%  ЦВЕТА
% ============================================================================
\definecolor{pastelblue}{RGB}{163,210,237}
\definecolor{pastelgreen}{RGB}{152,225,172}
\definecolor{pastelpurple}{RGB}{200,170,220}
\definecolor{pastelorange}{RGB}{255,210,170}
\definecolor{pastelteal}{RGB}{170,215,220}
\definecolor{pastelred}{RGB}{250,180,185}
\definecolor{accentblue}{RGB}{70,130,180}
\definecolor{sectioncolor}{RGB}{50,50,50}
\definecolor{softgray}{RGB}{140,140,140}

% ============================================================================
%  ИНТЕРВАЛ
% ============================================================================
\linespread{1.15}

% ============================================================================
%  ФОРМАТ ЗАГОЛОВКОВ СЕКЦИЙ
% ============================================================================
\titleformat{\section}
  {\Large\bfseries\color{sectioncolor}}
  {\thesection}{1em}{}
  [\vspace{-0.5em}{\color{pastelpurple}\hrule height 1.5pt}]

\titleformat{\subsection}
  {\large\bfseries\color{sectioncolor!85}}
  {\thesubsection}{0.8em}{}

\titleformat{\subsubsection}
  {\normalsize\bfseries\color{sectioncolor!70}}
  {\thesubsubsection}{0.6em}{}

% ============================================================================
%  КОЛОНТИТУЛЫ
% ============================================================================
\pagestyle{fancy}
\fancyhf{}
\fancyhead[L]{\small\color{softgray} MusicTech --- Центральный Университет}
\fancyhead[R]{\small\color{softgray} Score-Following 2026}
\fancyfoot[C]{\small\color{softgray}\thepage}
\renewcommand{\headrulewidth}{0.4pt}
\renewcommand{\headrule}{\hbox to\headwidth{\color{pastelpurple}\leaders\hrule height \headrulewidth\hfill}}

% ============================================================================
%  TCOLORBOX ОКРУЖЕНИЯ
% ============================================================================
\newtcolorbox{defbox}[1][]{
  enhanced, breakable,
  colback=pastelblue!12,
  colframe=pastelblue!60!black,
  sharp corners,
  boxrule=0.8pt,
  fonttitle=\bfseries,
  attach boxed title to top left={yshift=-2mm, xshift=5mm},
  boxed title style={
    colback=pastelblue!60!black,
    colframe=pastelblue!60!black,
    sharp corners
  },
  title={Определение},
  #1
}

\newtcolorbox{algbox}[1][]{
  enhanced, breakable,
  colback=pastelgreen!12,
  colframe=pastelgreen!60!black,
  sharp corners,
  boxrule=0.8pt,
  fonttitle=\bfseries,
  attach boxed title to top left={yshift=-2mm, xshift=5mm},
  boxed title style={
    colback=pastelgreen!60!black,
    colframe=pastelgreen!60!black,
    sharp corners
  },
  title={Алгоритм},
  #1
}

\newtcolorbox{exbox}[1][]{
  enhanced, breakable,
  colback=pastelorange!12,
  colframe=pastelorange!60!black,
  sharp corners,
  boxrule=0.8pt,
  fonttitle=\bfseries,
  attach boxed title to top left={yshift=-2mm, xshift=5mm},
  boxed title style={
    colback=pastelorange!60!black,
    colframe=pastelorange!60!black,
    sharp corners
  },
  title={Пример},
  #1
}

\newtcolorbox{warnbox}[1][]{
  enhanced, breakable,
  colback=pastelred!12,
  colframe=pastelred!60!black,
  sharp corners,
  boxrule=0.8pt,
  fonttitle=\bfseries,
  attach boxed title to top left={yshift=-2mm, xshift=5mm},
  boxed title style={
    colback=pastelred!60!black,
    colframe=pastelred!60!black,
    sharp corners
  },
  title={Внимание},
  #1
}

\newtcolorbox{notebox}[1][]{
  enhanced, breakable,
  colback=pastelteal!12,
  colframe=pastelteal!60!black,
  sharp corners,
  boxrule=0.8pt,
  fonttitle=\bfseries,
  attach boxed title to top left={yshift=-2mm, xshift=5mm},
  boxed title style={
    colback=pastelteal!60!black,
    colframe=pastelteal!60!black,
    sharp corners
  },
  title={Заметка},
  #1
}

\newtcolorbox{formulabox}[1][]{
  enhanced, breakable,
  colback=pastelpurple!12,
  colframe=pastelpurple!60!black,
  sharp corners,
  boxrule=0.8pt,
  fonttitle=\bfseries,
  attach boxed title to top left={yshift=-2mm, xshift=5mm},
  boxed title style={
    colback=pastelpurple!60!black,
    colframe=pastelpurple!60!black,
    sharp corners
  },
  title={Формула},
  #1
}

% ============================================================================
%  INLINE ЦВЕТНЫЕ КОМАНДЫ
% ============================================================================
\newcommand{\cblue}[1]{\textcolor{accentblue}{#1}}
\newcommand{\cgreen}[1]{\textcolor{pastelgreen!60!black}{#1}}
\newcommand{\cpurple}[1]{\textcolor{pastelpurple!70!black}{#1}}
\newcommand{\corange}[1]{\textcolor{pastelorange!70!black}{#1}}
\newcommand{\cteal}[1]{\textcolor{pastelteal!70!black}{#1}}

% --- Highlight команды ---
\newcommand{\hlblue}[1]{\colorbox{pastelblue!25}{#1}}
\newcommand{\hlgreen}[1]{\colorbox{pastelgreen!25}{#1}}
\newcommand{\hlpurple}[1]{\colorbox{pastelpurple!25}{#1}}
\newcommand{\hlorange}[1]{\colorbox{pastelorange!25}{#1}}
\newcommand{\hlteal}[1]{\colorbox{pastelteal!25}{#1}}

% ============================================================================
%  LISTINGS
% ============================================================================
\lstset{
  basicstyle=\ttfamily\small,
  keywordstyle=\color{accentblue}\bfseries,
  commentstyle=\color{softgray}\itshape,
  stringstyle=\color{pastelgreen!60!black},
  numbers=left,
  numberstyle=\tiny\color{softgray},
  numbersep=6pt,
  frame=single,
  framerule=0.4pt,
  rulecolor=\color{pastelpurple!40},
  backgroundcolor=\color{pastelpurple!5},
  breaklines=true,
  showstringspaces=false,
  tabsize=4,
  captionpos=b
}

\lstdefinestyle{python}{
  language=Python,
  morekeywords={self,None,True,False,as,with,yield,async,await}
}

% ============================================================================
%  ГРАФИКА
% ============================================================================
\graphicspath{{../figures/}}

% ============================================================================
\begin{document}

% ============================================================================
%  ТИТУЛЬНАЯ СТРАНИЦА
% ============================================================================
\begin{titlepage}
\begin{tikzpicture}[remember picture, overlay]
  % фоновые декоративные элементы
  \fill[pastelblue!20] (current page.south west) rectangle (current page.north east);

  \fill[pastelblue!35, opacity=0.6]
    ([xshift=-2cm, yshift=-2cm]current page.north east)
    circle (6cm);
  \fill[pastelpurple!30, opacity=0.5]
    ([xshift=3cm, yshift=3cm]current page.south west)
    circle (5cm);
  \fill[pastelgreen!25, opacity=0.5]
    ([xshift=-4cm, yshift=5cm]current page.south east)
    circle (3.5cm);
  \fill[pastelorange!20, opacity=0.4]
    ([xshift=6cm, yshift=-3cm]current page.north west)
    circle (4cm);
  \fill[pastelteal!25, opacity=0.4]
    ([xshift=0cm, yshift=-6cm]current page.north)
    circle (2.5cm);

  % горизонтальные декоративные полосы
  \draw[pastelpurple!60, line width=2pt]
    ([yshift=-4cm]current page.north west) -- ([yshift=-4cm]current page.north east);
  \draw[pastelblue!60, line width=1.5pt]
    ([yshift=-4.3cm]current page.north west) -- ([yshift=-4.3cm]current page.north east);

  \draw[pastelpurple!60, line width=2pt]
    ([yshift=3.5cm]current page.south west) -- ([yshift=3.5cm]current page.south east);
  \draw[pastelblue!60, line width=1.5pt]
    ([yshift=3.2cm]current page.south west) -- ([yshift=3.2cm]current page.south east);

  % декоративные нотные элементы
  \node[opacity=0.08, font=\fontsize{80}{80}\selectfont, rotate=15, color=pastelpurple!50]
    at ([xshift=5cm, yshift=-6cm]current page.north west) {$\sharp$};
  \node[opacity=0.06, font=\fontsize{60}{60}\selectfont, rotate=-10, color=pastelblue!50]
    at ([xshift=-3cm, yshift=6cm]current page.south east) {$\flat$};

  % основной текстовый блок
  \node[anchor=center] at ([yshift=5cm]current page.center) {
    \color{softgray}\large Мастерская MusicTech --- Направление 1: Логика работы
  };
  \node[anchor=center] at ([yshift=3.5cm]current page.center) {
    \color{sectioncolor}\fontsize{28}{34}\selectfont\bfseries Score-Following
  };
  \node[anchor=center, text width=12cm, align=center] at ([yshift=2.2cm]current page.center) {
    \color{accentblue}\fontsize{16}{20}\selectfont
    Алгоритмы отслеживания партитуры в реальном времени
  };
  \node[anchor=center] at ([yshift=0.8cm]current page.center) {
    \color{pastelpurple!70!black}\rule{8cm}{0.8pt}
  };
  \node[anchor=center] at ([yshift=-0.3cm]current page.center) {
    \color{sectioncolor}\Large Никита Новицкий
  };
  \node[anchor=center] at ([yshift=-1.2cm]current page.center) {
    \color{softgray}\large Центральный Университет (Т-Банк)
  };
  \node[anchor=center] at ([yshift=-2.8cm]current page.center) {
    \color{accentblue}\large 2026
  };

  % маленькие цветные кружки-декорации
  \foreach \x/\y/\c/\r in {
    -7/8/pastelblue/0.15,
    -6.5/7.5/pastelpurple/0.1,
    7/7/pastelgreen/0.12,
    6.5/6.5/pastelorange/0.08,
    -5/-7/pastelteal/0.13,
    6/-6/pastelred/0.1,
    0/9/pastelblue/0.08,
    -3/8.5/pastelpurple/0.06
  } {
    \fill[\c!80] ([xshift=\x cm, yshift=\y cm]current page.center) circle (\r cm);
  }
\end{tikzpicture}
\end{titlepage}

% ============================================================================
%  ОГЛАВЛЕНИЕ
% ============================================================================
\newpage
\tableofcontents
\newpage

% ============================================================================
%  РАЗДЕЛ 1: ВВЕДЕНИЕ
% ============================================================================
\section{Введение}
\label{sec:intro}

\subsection{Что такое score-following}

\textbf{Score-following} (отслеживание партитуры) --- это задача автоматического
определения текущей позиции живого исполнителя в нотном тексте (партитуре)
в реальном времени. Система принимает на вход поток музыкальных событий ---
MIDI-ноты или аудиосигнал --- и выдаёт соответствие между тем, что играет
исполнитель, и тем, что написано в нотах.

Живое исполнение никогда не совпадает с партитурой точно: музыкант допускает
отклонения в темпе (\emph{rubato}), динамике, может пропускать ноты или
добавлять орнаменты. Задача алгоритма --- быть устойчивым к этим отклонениям
и \emph{в реальном времени} указывать, какую ноту исполнитель играет прямо сейчас.

\subsection{Области применения}

\begin{enumerate}[leftmargin=2em]
  \item \textbf{Автоматический аккомпанемент.}
    Компьютер <<следит>> за солистом и подстраивает оркестровую партию
    под его темп. Первые такие системы появились в 1984~году.
  \item \textbf{Автоматический переворот страниц.}
    Электронный пульт (например, iPad с нотами) автоматически
    перелистывает страницу, когда исполнитель доходит до конца.
  \item \textbf{Визуализация в реальном времени.}
    Подсветка текущей ноты в партитуре для зрителей, синхронизация
    световых эффектов с музыкой.
  \item \textbf{Музыкальное образование.}
    Оценка точности исполнения ученика: где он ошибся, где отстал
    от темпа, какие ноты пропустил.
  \item \textbf{Интерактивные инсталляции.}
    Генеративная графика и звук, реагирующие на живое исполнение.
\end{enumerate}

\subsection{Краткая история}

Задача score-following была поставлена независимо двумя исследователями в 1984 году:
\cblue{Роджер Данненберг} (Roger Dannenberg, Carnegie Mellon)~\cite{dannenberg1984}
и \cblue{Барри Версо} (Barry Vercoe, MIT)~\cite{vercoe1984}. Оба предложили
алгоритмы на основе динамического программирования для отслеживания MIDI-клавиатуры.

В 1990--2000-х годах доминировали подходы на основе \emph{скрытых марковских
моделей} (HMM): Orio \& Déchelle~\cite{orio2001}, Raphael~\cite{cont2006}.
Параллельно развивались методы \emph{Dynamic Time Warping}
(DTW): Dixon~\cite{dixon2005}, Müller~\cite{muller2007}.

С 2018 года начали появляться нейросетевые подходы:
Dorfer et al.~\cite{dorfer2018} применили свёрточные сети для audio-to-score
alignment, а в 2025 году Henkel et al.~\cite{henkel2025} представили
HeurMiT --- нейросетевой фреймворк для полифонических инструментов.

\begin{defbox}[title={Постановка задачи score-following}]
\textbf{Дано:}
\begin{itemize}[leftmargin=1.5em]
  \item Партитура $S = \{(p_i,\; t_i^{\mathrm{on}},\; t_i^{\mathrm{off}})\}_{i=1}^{N}$
        --- последовательность нот с высотами $p_i$ и временными метками.
  \item Поток событий исполнения
        $P = \{(q_j,\; \tau_j^{\mathrm{on}},\; \tau_j^{\mathrm{off}})\}_{j=1}^{M}$,
        поступающий в реальном времени.
\end{itemize}
\textbf{Найти:}
функцию выравнивания $\varphi: \{1,\ldots,M\} \to \{1,\ldots,N\}$,
такую что $\varphi(j)$ --- позиция в партитуре, соответствующая
$j$-му событию исполнения.

\textbf{Требования:}
\begin{itemize}[leftmargin=1.5em]
  \item \emph{Реальное время:} $\varphi(j)$ вычисляется, используя
        только $p_1, \ldots, p_j$ (без знания будущих событий).
  \item \emph{Монотонность:} если $j_1 < j_2$, то $\varphi(j_1) \leq \varphi(j_2)$.
  \item \emph{Робастность:} устойчивость к ошибкам, пропускам, rubato.
\end{itemize}
\end{defbox}

\subsection{Pipeline системы score-following}

\begin{center}
\begin{tikzpicture}[
  box/.style={
    draw, rounded corners=3pt, minimum height=1.1cm, minimum width=2.8cm,
    font=\small\bfseries, text=sectioncolor, text width=2.4cm, align=center,
    line width=0.6pt
  },
  arr/.style={-{Stealth[length=6pt]}, line width=0.8pt, color=accentblue!80}
]
  \node[box, fill=pastelblue!30, draw=pastelblue!60!black] (midi) at (0,0) {MIDI вход};
  \node[box, fill=pastelgreen!30, draw=pastelgreen!60!black] (feat) at (3.6,0) {Извлечение признаков};
  \node[box, fill=pastelpurple!30, draw=pastelpurple!60!black] (algo) at (7.5,0) {Алгоритм выравнивания};
  \node[box, fill=pastelorange!30, draw=pastelorange!60!black] (pos)  at (11.4,0) {Позиция в партитуре};
  \node[box, fill=pastelteal!30, draw=pastelteal!60!black] (out)  at (15,0) {Аккомпа\-немент};

  \draw[arr] (midi) -- (feat);
  \draw[arr] (feat) -- (algo);
  \draw[arr] (algo) -- (pos);
  \draw[arr] (pos)  -- (out);

  \node[above=0.15cm of midi, font=\scriptsize\color{softgray}] {Клавиатура};
  \node[above=0.15cm of feat, font=\scriptsize\color{softgray}] {Pitch, velocity};
  \node[above=0.15cm of algo, font=\scriptsize\color{softgray}] {DTW / HMM / OLTW};
  \node[above=0.15cm of pos, font=\scriptsize\color{softgray}] {Нота $\varphi(j)$};
  \node[above=0.15cm of out, font=\scriptsize\color{softgray}] {Синтезатор};
\end{tikzpicture}
\end{center}


% ============================================================================
%  РАЗДЕЛ 2: MIDI И МУЗЫКАЛЬНЫЕ ОСНОВЫ
% ============================================================================
\section{MIDI и музыкальные основы}
\label{sec:midi}

\subsection{Протокол MIDI}

MIDI (\emph{Musical Instrument Digital Interface}) --- это стандарт цифровой
передачи музыкальных данных, разработанный в 1983~году. В отличие от аудиосигнала,
MIDI передаёт не звуковую волну, а \emph{символические события}: какую ноту нажали,
с какой силой, когда отпустили.

Основные типы сообщений, релевантные для score-following:

\begin{itemize}[leftmargin=2em]
  \item \hlblue{\texttt{Note On} (\texttt{0x90})} --- нажатие клавиши.
        Содержит номер ноты (pitch) и силу нажатия (velocity).
  \item \hlblue{\texttt{Note Off} (\texttt{0x80})} --- отпускание клавиши.
        Содержит номер ноты и velocity отпускания (обычно~0).
  \item \texttt{Control Change} (\texttt{0xB0}) --- педаль, модуляция и др.
  \item \texttt{Program Change} (\texttt{0xC0}) --- смена инструмента.
\end{itemize}

\begin{defbox}[title={MIDI-событие}]
Каждое MIDI-событие представляется кортежем:
\[
  e = (\textit{type},\; \textit{pitch},\; \textit{velocity},\; \textit{timestamp})
\]
где:
\begin{itemize}[leftmargin=1.5em]
  \item $\textit{type} \in \{\texttt{note\_on},\; \texttt{note\_off}\}$
  \item $\textit{pitch} \in \{0, 1, \ldots, 127\}$ --- номер ноты по стандарту MIDI
  \item $\textit{velocity} \in \{0, 1, \ldots, 127\}$ --- сила нажатия
  \item $\textit{timestamp} \in \mathbb{R}_{\geq 0}$ --- время в секундах
\end{itemize}
\end{defbox}

\subsection{Таблица нот MIDI}

\begin{center}
\begin{tabular}{lccccccccccccc}
\toprule
\textbf{Нота} & C3 & D3 & E3 & F3 & G3 & A3 & B3 & C4 & D4 & E4 & F4 & G4 & A4 \\
\midrule
\textbf{MIDI} & 48 & 50 & 52 & 53 & 55 & 57 & 59 & 60 & 62 & 64 & 65 & 67 & 69 \\
\bottomrule
\end{tabular}
\captionof{table}{Соответствие нот и номеров MIDI. C4 (<<до>> первой октавы) = 60.}
\label{tab:midi_notes}
\end{center}

\subsection{Клавиатура и MIDI-номера}

\begin{center}
\begin{tikzpicture}[scale=0.85]
  % белые клавиши
  \foreach \i/\note/\midi in {
    0/C4/60, 1/D4/62, 2/E4/64, 3/F4/65, 4/G4/67, 5/A4/69, 6/B4/71
  } {
    \draw[fill=white, draw=sectioncolor!60, line width=0.5pt]
      (\i*1.4, 0) rectangle ++(1.3, 4);
    \node[font=\scriptsize\bfseries, color=sectioncolor] at (\i*1.4+0.65, 0.5) {\note};
    \node[font=\tiny, color=accentblue] at (\i*1.4+0.65, 1.2) {\midi};
  }
  % чёрные клавиши
  \foreach \i/\note/\midi in {
    0/C\#/61, 1/D\#/63, 3/F\#/66, 4/G\#/68, 5/A\#/70
  } {
    \draw[fill=sectioncolor!85, draw=sectioncolor, line width=0.5pt]
      (\i*1.4+0.85, 1.8) rectangle ++(0.9, 2.2);
    \node[font=\tiny\bfseries, color=white] at (\i*1.4+1.3, 2.5) {\midi};
  }
\end{tikzpicture}
\end{center}

\subsection{Представление партитуры и исполнения}

Для алгоритмов score-following нам нужны два формальных представления:

\begin{defbox}[title={Партитура (Score)}]
Последовательность нотных событий:
\[
  S = \bigl\{(p_i,\; t_i^{\mathrm{on}},\; t_i^{\mathrm{off}})\bigr\}_{i=1}^{N}
\]
где $p_i \in \{0,\ldots,127\}$ --- высота $i$-й ноты, $t_i^{\mathrm{on}}$
и $t_i^{\mathrm{off}}$ --- номинальные времена начала и конца ноты по партитуре.
$N$ --- общее количество нот в произведении.
\end{defbox}

\begin{defbox}[title={Исполнение (Performance)}]
Последовательность реальных событий от исполнителя:
\[
  P = \bigl\{(q_j,\; \tau_j^{\mathrm{on}},\; \tau_j^{\mathrm{off}})\bigr\}_{j=1}^{M}
\]
где $q_j$ --- реально сыгранная высота, $\tau_j^{\mathrm{on}}$ и $\tau_j^{\mathrm{off}}$
--- реальные времена. $M$ может отличаться от $N$: исполнитель может пропускать
или добавлять ноты.
\end{defbox}


% ============================================================================
%  РАЗДЕЛ 3: ФОРМАЛИЗАЦИЯ ЗАДАЧИ
% ============================================================================
\section{Формализация задачи}
\label{sec:formalization}

\subsection{Функция выравнивания}

\begin{formulabox}[title={Функция выравнивания}]
Задача score-following --- построить отображение:
\[
  \varphi: \{1, \ldots, M\} \longrightarrow \{1, \ldots, N\}
\]
которое сопоставляет каждому событию исполнения $j$ позицию
в партитуре $\varphi(j)$.
\end{formulabox}

\subsection{Ограничение монотонности}

Ключевое ограничение: исполнитель играет партитуру <<вперёд>>
(с возможными пропусками, но без возврата назад):
\[
  j_1 < j_2 \implies \varphi(j_1) \leq \varphi(j_2)
\]

Равенство $\varphi(j_1) = \varphi(j_2)$ допускается: несколько
событий исполнения могут относиться к одной ноте партитуры
(например, повторные нажатия, трели, орнаменты).

\subsection{Функция стоимости}

\begin{formulabox}[title={Функция стоимости (cost function)}]
Для сравнения ноты партитуры $s_i$ и события исполнения $p_j$
определим локальную стоимость:
\[
  d(s_i, p_j) = |p_i - q_j|
\]
где $p_i$ --- pitch ноты в партитуре, $q_j$ --- pitch сыгранной ноты.

Если pitch совпадает: $d = 0$. Полутоновая ошибка: $d = 1$.
Октавная ошибка: $d = 12$.
\end{formulabox}

\subsection{Оптимальное выравнивание}

Оптимальное выравнивание минимизирует суммарную стоимость:

\begin{formulabox}[title={Оптимальное выравнивание}]
\[
  \varphi^* = \arg\min_{\varphi} \sum_{j=1}^{M} d\bigl(s_{\varphi(j)},\; p_j\bigr)
\]
при ограничении монотонности $\varphi$.
\end{formulabox}

\subsection{Online vs Offline}

Принципиальное различие двух постановок:

\begin{itemize}[leftmargin=2em]
  \item \textbf{Offline:} доступно всё исполнение $P$ целиком.
        Можно вычислить глобально оптимальное $\varphi^*$.
        Применение: анализ записей, разметка данных.
  \item \textbf{Online:} в момент $j$ доступны только $p_1, \ldots, p_j$.
        Нужно выдать $\varphi(j)$ немедленно, без знания будущих событий.
        Применение: живой аккомпанемент, переворот страниц.
\end{itemize}

В проекте MusicTech нас интересует именно \emph{online}-постановка.

\subsection{Визуализация выравнивания}

\begin{center}
\begin{tikzpicture}[
  note/.style={circle, draw, minimum size=0.7cm, font=\small\bfseries, line width=0.5pt},
  arr/.style={-{Stealth[length=5pt]}, color=accentblue!70, line width=0.6pt, dashed}
]
  % партитура (верхний ряд)
  \node[font=\small\bfseries\color{softgray}] at (-1.5, 1.5) {Партитура:};
  \foreach \i/\p in {1/C4, 2/E4, 3/G4, 4/A4, 5/C5, 6/E5} {
    \node[note, fill=pastelblue!30, draw=pastelblue!60!black] (s\i) at (\i*1.6, 1.5) {\p};
  }

  % исполнение (нижний ряд)
  \node[font=\small\bfseries\color{softgray}] at (-1.5, -0.5) {Исполнение:};
  \foreach \j/\q in {1/C4, 2/E4, 3/E4, 4/G4, 5/A4, 6/B4, 7/C5, 8/E5} {
    \node[note, fill=pastelorange!30, draw=pastelorange!60!black] (p\j) at (\j*1.3-0.2, -0.5) {\q};
  }

  % стрелки выравнивания
  \draw[arr] (p1) -- (s1);
  \draw[arr] (p2) -- (s2);
  \draw[arr] (p3) -- (s2);
  \draw[arr] (p4) -- (s3);
  \draw[arr] (p5) -- (s4);
  \draw[arr] (p6) -- (s4);
  \draw[arr] (p7) -- (s5);
  \draw[arr] (p8) -- (s6);

  \node[font=\scriptsize\color{softgray}, anchor=north west, text width=3.5cm] at (11.5, 1)
    {$\varphi$: $1{\to}1$, $2{\to}2$, $3{\to}2$, $4{\to}3$, $5{\to}4$, $6{\to}4$, $7{\to}5$, $8{\to}6$};
\end{tikzpicture}
\end{center}

На диаграмме видно, что события~3 и~6 исполнения связываются с
нотами~2 и~4 партитуры соответственно (повторные нажатия).
Нота B4 (событие~6) выравнивается к A4 --- это ошибка исполнителя,
и алгоритм выбирает ближайшую ноту партитуры.


% ============================================================================
%  РАЗДЕЛ 4: DYNAMIC TIME WARPING
% ============================================================================
\section{Dynamic Time Warping (DTW)}
\label{sec:dtw}

\subsection{Определение}

\begin{defbox}[title={Dynamic Time Warping}]
DTW (\emph{Dynamic Time Warping}, динамическая трансформация временной шкалы)~\cite{sakoe1978}
--- это алгоритм поиска оптимального выравнивания двух временных
последовательностей, допускающий нелинейные деформации по времени.

Даны две последовательности:
\[
  X = (x_1, x_2, \ldots, x_N), \quad Y = (y_1, y_2, \ldots, y_M)
\]
DTW находит путь через матрицу стоимости, минимизирующий
суммарное расстояние между выровненными элементами.
\end{defbox}

\subsection{Матрица локальной стоимости}

Для каждой пары $(i,j)$ вычисляем локальную стоимость:
\[
  C(i,j) = d(x_i, y_j)
\]
В случае MIDI-нот: $d(x_i, y_j) = |p_i - q_j|$ (разница pitch).
Для audio: евклидово расстояние между спектральными векторами.

Матрица $C$ имеет размер $N \times M$.

\subsection{Рекурсия DTW}

\begin{formulabox}[title={Рекурсия DTW}]
Матрица накопленной стоимости $D$ вычисляется рекурсивно:
\[
  D(i,j) = C(i,j) + \min
  \begin{cases}
    D(i-1,\; j)   & \text{(вставка)} \\
    D(i,\; j-1)   & \text{(удаление)} \\
    D(i-1,\; j-1) & \text{(совпадение)}
  \end{cases}
\]
с граничными условиями:
\[
  D(0,0) = 0, \quad D(i,0) = \infty \;\forall\, i > 0, \quad D(0,j) = \infty \;\forall\, j > 0
\]
\end{formulabox}

Три варианта перехода имеют интуитивный смысл:
\begin{itemize}[leftmargin=2em]
  \item $D(i-1, j-1) \to D(i,j)$: диагональный шаг --- следующая нота
        партитуры сопоставляется со следующей нотой исполнения.
  \item $D(i-1, j) \to D(i,j)$: вертикальный шаг --- нота партитуры
        пропущена исполнителем.
  \item $D(i, j-1) \to D(i,j)$: горизонтальный шаг --- исполнитель
        сыграл лишнюю ноту (орнамент, ошибка).
\end{itemize}

\subsection{Оптимальный путь}

Итоговая стоимость выравнивания: $D(N, M)$.
Оптимальный путь $W = (w_1, w_2, \ldots, w_K)$, где $w_k = (i_k, j_k)$,
находится обратным проходом (\emph{backtracking}) от $(N,M)$ к $(1,1)$:
на каждом шаге выбираем предшественника с минимальным значением $D$.

\subsection{Псевдокод DTW}

\begin{algbox}[title={Алгоритм DTW}]
\begin{lstlisting}[style=python, numbers=none]
def dtw(X, Y):
    N, M = len(X), len(Y)
    D = matrix(N+1, M+1, fill=inf)
    D[0][0] = 0

    for i in range(1, N+1):
        for j in range(1, M+1):
            cost = abs(X[i-1] - Y[j-1])
            D[i][j] = cost + min(
                D[i-1][j],      # insertion
                D[i][j-1],      # deletion
                D[i-1][j-1]     # match
            )

    # backtracking
    path = [(N, M)]
    i, j = N, M
    while i > 1 or j > 1:
        candidates = []
        if i > 1: candidates.append((D[i-1][j], i-1, j))
        if j > 1: candidates.append((D[i][j-1], i, j-1))
        if i > 1 and j > 1:
            candidates.append((D[i-1][j-1], i-1, j-1))
        _, i, j = min(candidates)
        path.append((i, j))

    return D[N][M], reversed(path)
\end{lstlisting}
\end{algbox}

\subsection{Сложность}

\begin{formulabox}[title={Сложность DTW}]
\begin{itemize}[leftmargin=1.5em]
  \item \textbf{Время:} $O(N \cdot M)$
  \item \textbf{Память:} $O(N \cdot M)$
\end{itemize}
Для произведения из $N = 1000$ нот и исполнения из $M = 1100$ нот:
$\approx 1{,}1 \cdot 10^6$ операций --- мгновенно для offline.
\end{formulabox}

\subsection{Ограничение Sakoe--Chiba}

Для ускорения DTW можно ограничить область вычислений полосой
шириной $R$ вдоль диагонали~\cite{sakoe1978}:
\[
  |i - j| \leq R
\]
Это даёт сложность $O(N \cdot R)$ вместо $O(N \cdot M)$. Параметр $R$
выбирается исходя из максимального ожидаемого отклонения темпа.

\subsection{Визуализация: матрица стоимости и оптимальный путь}

\begin{center}
\begin{tikzpicture}[scale=0.75]
  % матрица 6x6
  \def\data{
    {0, 2, 5, 7, 12, 12},
    {2, 0, 3, 5, 10, 10},
    {4, 2, 0, 2, 7,  7},
    {5, 3, 2, 0, 5,  5},
    {7, 5, 3, 2, 0,  0},
    {12,10,7, 5, 0,  0}
  }

  % нарисовать ячейки с цветом по значению
  \foreach \i in {0,...,5} {
    \foreach \j in {0,...,5} {
      \pgfmathsetmacro{\val}{{\data}[\i][\j]}
      \pgfmathsetmacro{\intensity}{min(100, \val*7)}
      \fill[pastelblue!\intensity!white] (\j, 5-\i) rectangle ++(1,1);
      \draw[softgray!40] (\j, 5-\i) rectangle ++(1,1);
      \node[font=\tiny] at (\j+0.5, 5.5-\i) {\pgfmathprintnumber[fixed]{\val}};
    }
  }

  % оптимальный путь
  \draw[red!70!black, line width=1.8pt, -{Stealth[length=5pt]}]
    (0.5, 5.5) -- (1.5, 4.5) -- (2.5, 3.5) -- (3.5, 2.5) -- (4.5, 1.5) -- (5.5, 0.5);

  % подписи
  \node[font=\small\bfseries, color=sectioncolor] at (3, -0.5) {Исполнение $Y$};
  \node[font=\small\bfseries, color=sectioncolor, rotate=90] at (-0.7, 3) {Партитура $X$};

  % номера
  \foreach \j/\note in {0/C4, 1/D4, 2/F4, 3/G4, 4/C5, 5/C5} {
    \node[font=\tiny\color{accentblue}] at (\j+0.5, 6.3) {\note};
  }
  \foreach \i/\note in {0/C4, 1/D4, 2/F4, 3/G4, 4/C5, 5/C5} {
    \node[font=\tiny\color{accentblue}] at (-0.3, 5.5-\i) {\note};
  }
\end{tikzpicture}
\end{center}

Красная линия --- оптимальный путь. Тёмные ячейки --- высокая стоимость,
светлые --- низкая. Путь проходит по <<долине>> минимальных значений.

\subsection{Числовой пример}

\begin{exbox}[title={DTW: числовой пример}]
Пусть партитура содержит 5 нот: $X = (60, 64, 67, 72, 76)$ (C4, E4, G4, C5, E5).

Исполнитель сыграл: $Y = (60, 63, 67, 67, 72, 76)$ (ошибка во второй ноте:
63 вместо 64, повторение G4).

\medskip
\textbf{Шаг 1.} Матрица локальной стоимости $C(i,j) = |x_i - y_j|$:

\begin{center}
\begin{tabular}{c|cccccc}
  & $y_1{=}60$ & $y_2{=}63$ & $y_3{=}67$ & $y_4{=}67$ & $y_5{=}72$ & $y_6{=}76$ \\
\midrule
$x_1{=}60$ & 0 & 3 & 7 & 7 & 12 & 16 \\
$x_2{=}64$ & 4 & 1 & 3 & 3 & 8  & 12 \\
$x_3{=}67$ & 7 & 4 & 0 & 0 & 5  & 9 \\
$x_4{=}72$ & 12& 9 & 5 & 5 & 0  & 4 \\
$x_5{=}76$ & 16& 13& 9 & 9 & 4  & 0 \\
\end{tabular}
\end{center}

\textbf{Шаг 2.} Заполняем $D(i,j)$ по рекурсии:

$D(1,1) = 0$,\quad $D(1,2) = 0 + 3 = 3$,\quad $D(2,2) = 1 + \min(4, 3, 0) = 1$\quad и т.д.

\textbf{Шаг 3.} Оптимальный путь: $(1,1) \to (2,2) \to (3,3) \to (3,4) \to (4,5) \to (5,6)$.

Итоговая стоимость: $D(5,6) = 0 + 1 + 0 + 0 + 0 + 0 = 1$.

Выравнивание: $\varphi = (1, 2, 3, 3, 4, 5)$ --- третья и четвёртая ноты
исполнения обе привязаны к $x_3 = \text{G4}$.
\end{exbox}

\begin{warnbox}[title={Ограничения DTW}]
\begin{enumerate}[leftmargin=1.5em]
  \item \textbf{Offline-алгоритм:} требует полного исполнения для вычисления.
        Не применим для real-time score-following напрямую.
  \item \textbf{Квадратичная сложность:} $O(NM)$ может быть слишком медленно
        для длинных произведений в реальном времени.
  \item \textbf{Нет вероятностной модели:} DTW не моделирует неопределённость
        исполнителя, не даёт распределения вероятностей по позициям.
\end{enumerate}
\end{warnbox}


% ============================================================================
%  РАЗДЕЛ 5: ONLINE TIME WARPING
% ============================================================================
\section{Online Time Warping (OLTW)}
\label{sec:oltw}

\subsection{Определение}

\begin{defbox}[title={Online Time Warping}]
OLTW (\emph{Online Time Warping})~\cite{dixon2005} --- инкрементальная
версия DTW, обрабатывающая исполнение событие за событием.

При поступлении нового наблюдения $y_j$ алгоритм:
\begin{enumerate}[leftmargin=1.5em]
  \item Вычисляет новый столбец матрицы $D$.
  \item Продлевает оптимальный путь на один шаг.
  \item Выдаёт текущую позицию в партитуре $\varphi(j)$.
\end{enumerate}
\end{defbox}

\subsection{Инкрементальное обновление}

При поступлении $y_j$ вычисляем $j$-й столбец $D$:

\begin{formulabox}[title={OLTW: обновление столбца}]
Для каждого $i = 1, \ldots, N$:
\[
  D(i,j) = C(i,j) + \min\bigl\{D(i-1,j),\; D(i,j-1),\; D(i-1,j-1)\bigr\}
\]
где $C(i,j) = |x_i - y_j|$.

Столбец $D(\cdot, j-1)$ уже вычислен на предыдущем шаге.
\end{formulabox}

\subsection{Механизм RunCount}

Ключевая проблема online-выравнивания: алгоритм может <<застрять>>
на одной ноте партитуры (бесконечно продлевая горизонтальный шаг)
или наоборот <<убежать>> вперёд. Dixon~\cite{dixon2005} предложил
механизм \textbf{RunCount}:

\begin{formulabox}[title={RunCount}]
$\mathrm{RunCount}$ --- счётчик последовательных шагов в одном направлении.
\begin{itemize}[leftmargin=1.5em]
  \item Если текущий шаг в том же направлении, что предыдущий:
        $\mathrm{RunCount} \gets \mathrm{RunCount} + 1$
  \item Иначе: $\mathrm{RunCount} \gets 0$
\end{itemize}
Если $\mathrm{RunCount} > \mathrm{MaxRun}$, принудительно выбираем
диагональный шаг.

Типичное значение: $\mathrm{MaxRun} = 3\text{--}5$.
\end{formulabox}

\subsection{Псевдокод OLTW}

\begin{algbox}[title={Алгоритм OLTW}]
\begin{lstlisting}[style=python, numbers=none]
class OLTW:
    def __init__(self, score, max_run=3):
        self.score = score
        self.N = len(score)
        self.max_run = max_run
        self.position = 0
        self.run_count = 0
        self.prev_direction = None
        self.prev_col = [inf] * (self.N + 1)
        self.prev_col[0] = 0

    def process(self, observation):
        curr_col = [inf] * (self.N + 1)
        backptr = [None] * (self.N + 1)

        for i in range(1, self.N + 1):
            cost = abs(self.score[i-1] - observation)
            options = [
                (self.prev_col[i],   "horiz"),
                (curr_col[i-1],      "vert"),
                (self.prev_col[i-1], "diag"),
            ]
            best_val, best_dir = min(options)
            curr_col[i] = cost + best_val
            backptr[i] = best_dir

        direction = backptr[self.position + 1]
        if direction == self.prev_direction:
            self.run_count += 1
        else:
            self.run_count = 0

        if self.run_count > self.max_run:
            direction = "diag"
            self.run_count = 0

        if direction in ("diag", "vert"):
            self.position = min(self.position + 1,
                                self.N - 1)
        self.prev_direction = direction
        self.prev_col = curr_col
        return self.position
\end{lstlisting}
\end{algbox}

\subsection{Сложность}

\begin{itemize}[leftmargin=2em]
  \item \textbf{На каждое новое наблюдение:} $O(N)$ --- один проход по столбцу.
  \item \textbf{Память:} $O(N)$ --- храним только два столбца (текущий и предыдущий).
  \item \textbf{Общая сложность:} $O(NM)$, но распределённая во времени:
        каждый шаг занимает $O(N)$, что приемлемо для real-time при $N < 10\,000$.
\end{itemize}

\subsection{Визуализация: пошаговое построение пути}

\begin{center}
\begin{tikzpicture}[scale=0.55]
  % Panel 1: t=3
  \begin{scope}[xshift=0cm]
    \draw[softgray!30] (0,0) grid (3,6);
    \draw[sectioncolor!50, line width=0.5pt] (0,0) rectangle (3,6);
    \draw[pastelgreen!70!black, line width=2pt, -{Stealth[length=5pt]}]
      (0.5,5.5) -- (1.5,4.5) -- (2.5,3.5);
    \fill[pastelgreen!40] (0.5,5.5) circle (0.2);
    \fill[pastelgreen!40] (1.5,4.5) circle (0.2);
    \fill[pastelgreen!60] (2.5,3.5) circle (0.2);
    \node[font=\small\bfseries, color=sectioncolor] at (1.5, -0.7) {$t = 3$};
    \node[font=\scriptsize, color=softgray] at (1.5, -1.3) {3 события};
  \end{scope}

  % Panel 2: t=6
  \begin{scope}[xshift=5cm]
    \draw[softgray!30] (0,0) grid (6,6);
    \draw[sectioncolor!50, line width=0.5pt] (0,0) rectangle (6,6);
    \draw[pastelblue!70!black, line width=2pt, -{Stealth[length=5pt]}]
      (0.5,5.5) -- (1.5,4.5) -- (2.5,3.5) -- (3.5,3.5) -- (4.5,2.5) -- (5.5,1.5);
    \fill[pastelblue!40] (0.5,5.5) circle (0.2);
    \fill[pastelblue!40] (1.5,4.5) circle (0.2);
    \fill[pastelblue!40] (2.5,3.5) circle (0.2);
    \fill[pastelblue!40] (3.5,3.5) circle (0.2);
    \fill[pastelblue!40] (4.5,2.5) circle (0.2);
    \fill[pastelblue!60] (5.5,1.5) circle (0.2);
    \node[font=\small\bfseries, color=sectioncolor] at (3, -0.7) {$t = 6$};
    \node[font=\scriptsize, color=softgray] at (3, -1.3) {6 событий};
  \end{scope}

  % Panel 3: t=9
  \begin{scope}[xshift=13cm]
    \draw[softgray!30] (0,0) grid (9,6);
    \draw[sectioncolor!50, line width=0.5pt] (0,0) rectangle (9,6);
    \draw[pastelpurple!70!black, line width=2pt, -{Stealth[length=5pt]}]
      (0.5,5.5) -- (1.5,4.5) -- (2.5,3.5) -- (3.5,3.5) -- (4.5,2.5)
      -- (5.5,1.5) -- (6.5,1.5) -- (7.5,0.5) -- (8.5,0.5);
    \foreach \x/\y in {0.5/5.5, 1.5/4.5, 2.5/3.5, 3.5/3.5, 4.5/2.5,
                        5.5/1.5, 6.5/1.5, 7.5/0.5} {
      \fill[pastelpurple!40] (\x,\y) circle (0.2);
    }
    \fill[pastelpurple!60] (8.5,0.5) circle (0.2);
    \node[font=\small\bfseries, color=sectioncolor] at (4.5, -0.7) {$t = 9$};
    \node[font=\scriptsize, color=softgray] at (4.5, -1.3) {9 событий};
  \end{scope}
\end{tikzpicture}
\end{center}

Три панели показывают рост пути OLTW по мере поступления новых событий.
Горизонтальные сегменты --- исполнитель задерживается на одной ноте,
диагональные --- синхронное продвижение.

\subsection{Сравнение DTW и OLTW}

\begin{center}
\begin{tcolorbox}[
  enhanced, breakable,
  colback=pastelteal!8,
  colframe=pastelteal!50!black,
  sharp corners,
  boxrule=0.6pt,
  title={\bfseries Сравнение DTW и OLTW},
  fonttitle=\bfseries,
  attach boxed title to top left={yshift=-2mm, xshift=5mm},
  boxed title style={colback=pastelteal!50!black, sharp corners}
]
\begin{tabularx}{\textwidth}{lXX}
\toprule
\textbf{Критерий} & \textbf{DTW} & \textbf{OLTW} \\
\midrule
Режим работы & Offline & Online (real-time) \\
Сложность (время) & $O(NM)$ & $O(N)$ на событие \\
Сложность (память) & $O(NM)$ & $O(N)$ \\
Оптимальность & Глобальный оптимум & Локально оптимальный \\
Устойчивость к rubato & Высокая & Средняя (RunCount) \\
Латентность & Нулевая (offline) & $< 10$\,мс \\
Применение & Анализ записей & Живой аккомпанемент \\
\bottomrule
\end{tabularx}
\end{tcolorbox}
\end{center}


% ============================================================================
%  РАЗДЕЛ 6: СКРЫТЫЕ МАРКОВСКИЕ МОДЕЛИ
% ============================================================================
\section{Скрытые марковские модели (HMM)}
\label{sec:hmm}

\subsection{HMM для score-following}

\begin{defbox}[title={HMM для score-following}]
Скрытая марковская модель (Hidden Markov Model, HMM)~\cite{rabiner1989}
моделирует score-following как задачу декодирования скрытых состояний:

\begin{itemize}[leftmargin=1.5em]
  \item \textbf{Скрытые состояния} $s_1, s_2, \ldots, s_N$ --- позиции
        в партитуре (каждая нота --- отдельное состояние).
  \item \textbf{Наблюдения} $o_1, o_2, \ldots, o_T$ --- MIDI-события
        от исполнителя.
  \item \textbf{Параметры модели} $\lambda = (A, B, \pi)$.
\end{itemize}
\end{defbox}

\subsection{Параметры модели}

\subsubsection{Матрица переходов $A$}

Матрица $A = \{a_{ij}\}$ задаёт вероятности перехода между состояниями:
\[
  a_{ij} = P(s_t = j \mid s_{t-1} = i)
\]

Для score-following матрица имеет ленточную структуру:

\begin{formulabox}[title={Матрица переходов}]
\[
  a_{ij} =
  \begin{cases}
    p_{\text{stay}}    & \text{если } j = i \quad\quad\;\;\, \text{(остаёмся на ноте)} \\
    p_{\text{advance}} & \text{если } j = i + 1 \quad \text{(переход к следующей)} \\
    p_{\text{skip}}    & \text{если } j = i + 2 \quad \text{(пропуск одной ноты)} \\
    0                  & \text{иначе}
  \end{cases}
\]
Типичные значения: $p_{\text{stay}} \approx 0.3$, $p_{\text{advance}} \approx 0.6$,
$p_{\text{skip}} \approx 0.1$.
\end{formulabox}

Это отражает реальное поведение исполнителя:
\begin{itemize}[leftmargin=2em]
  \item С вероятностью 0.3 длится текущая нота (долгие ноты, фермата).
  \item С вероятностью 0.6 переход к следующей ноте.
  \item С вероятностью 0.1 пропуск ноты (ошибка, сокращение).
\end{itemize}

\subsubsection{Функция эмиссии $B$}

Вероятность наблюдения $o_t$ в состоянии $s_t = i$:

\begin{formulabox}[title={Функция эмиссии}]
\[
  b_i(o_t) = P(o_t \mid s_t = i) = \mathcal{N}\bigl(o_t;\; \text{pitch}_i,\; \sigma^2\bigr)
\]
где $\text{pitch}_i$ --- номинальная высота $i$-й ноты партитуры,
$\sigma$ --- стандартное отклонение (обычно $\sigma \approx 1\text{--}3$
полутона).

Гауссова эмиссия обеспечивает:
\begin{itemize}[leftmargin=1.5em]
  \item Максимальная вероятность при точном совпадении pitch.
  \item Плавное убывание при полутоновых ошибках.
  \item Малая вероятность при октавных ошибках ($|o_t - \text{pitch}_i| = 12$).
\end{itemize}
\end{formulabox}

\subsubsection{Начальное распределение $\pi$}

\[
  \pi_i = P(s_1 = i) =
  \begin{cases}
    1 & \text{если } i = 1 \\
    0 & \text{иначе}
  \end{cases}
\]
Исполнение начинается с первой ноты партитуры.

\subsection{Алгоритм Forward}

Алгоритм Forward вычисляет вероятность наблюдения последовательности
$o_1, \ldots, o_t$ при нахождении в состоянии $j$ в момент $t$:

\begin{formulabox}[title={Forward-переменные}]
\textbf{Инициализация:}
\[
  \alpha_1(i) = \pi_i \cdot b_i(o_1)
\]
\textbf{Рекурсия:} для $t = 2, \ldots, T$:
\[
  \alpha_t(j) = \left[\sum_{i=1}^{N} \alpha_{t-1}(i) \cdot a_{ij}\right] \cdot b_j(o_t)
\]
\textbf{Текущая позиция:}
\[
  s_t^* = \arg\max_{j} \alpha_t(j)
\]
\end{formulabox}

Forward-алгоритм можно использовать для online score-following:
при каждом новом наблюдении $o_t$ обновляем вектор $\alpha_t$
за $O(N^2)$ операций (или $O(N)$, учитывая ленточную структуру $A$).

\subsection{Алгоритм Viterbi}

Алгоритм Viterbi находит \emph{наиболее вероятную последовательность состояний}:

\begin{formulabox}[title={Алгоритм Viterbi}]
\textbf{Инициализация:}
\[
  \delta_1(i) = \pi_i \cdot b_i(o_1), \qquad \psi_1(i) = 0
\]
\textbf{Рекурсия:} для $t = 2, \ldots, T$:
\[
  \delta_t(j) = \max_{i} \bigl[\delta_{t-1}(i) \cdot a_{ij}\bigr] \cdot b_j(o_t)
\]
\[
  \psi_t(j) = \arg\max_{i} \bigl[\delta_{t-1}(i) \cdot a_{ij}\bigr]
\]
\textbf{Обратный проход (backtracking):}
\[
  s_T^* = \arg\max_{i} \delta_T(i), \qquad s_t^* = \psi_{t+1}(s_{t+1}^*) \quad \text{для } t = T{-}1, \ldots, 1
\]
\end{formulabox}

\subsection{Диаграмма состояний HMM}

\begin{center}
\begin{tikzpicture}[
  state/.style={
    circle, draw=pastelpurple!60!black, fill=pastelpurple!20,
    minimum size=1.2cm, font=\small\bfseries, line width=0.6pt
  },
  trans/.style={-{Stealth[length=5pt]}, line width=0.6pt},
  obs/.style={
    rectangle, draw=pastelorange!60!black, fill=pastelorange!15,
    minimum height=0.7cm, minimum width=0.9cm, font=\scriptsize, rounded corners=2pt
  },
  emit/.style={-{Stealth[length=4pt]}, line width=0.4pt, dashed, color=pastelorange!60!black}
]
  % состояния
  \foreach \i/\note in {1/C4, 2/E4, 3/G4, 4/A4, 5/C5} {
    \node[state] (s\i) at (\i*2.5, 2) {$s_\i$};
    \node[font=\tiny\color{accentblue}] at (\i*2.5, 2.9) {\note};
  }

  % переходы advance
  \foreach \i [evaluate=\i as \j using int(\i+1)] in {1,...,4} {
    \draw[trans, color=pastelgreen!60!black]
      (s\i) -- node[above, font=\tiny, color=pastelgreen!60!black] {0.6} (s\j);
  }

  % переходы skip
  \foreach \i [evaluate=\i as \j using int(\i+2)] in {1,...,3} {
    \draw[trans, color=pastelred!60!black, bend left=35]
      (s\i) to node[above, font=\tiny, color=pastelred!60!black] {0.1} (s\j);
  }

  % self-loops
  \foreach \i in {1,...,5} {
    \draw[trans, color=accentblue!70] (s\i) edge[loop above, looseness=6, min distance=10mm]
      node[above, font=\tiny, color=accentblue!70] {0.3} (s\i);
  }

  % наблюдения
  \foreach \i in {1,...,5} {
    \node[obs] (o\i) at (\i*2.5, -0.3) {$o_\i$};
    \draw[emit] (s\i) -- (o\i);
  }

  % легенда
  \node[font=\scriptsize, color=pastelgreen!60!black] at (12, 4)
    {$\longrightarrow$ $p_{\text{advance}} = 0.6$};
  \node[font=\scriptsize, color=accentblue!70] at (12, 3.5)
    {$\circlearrowleft$ $p_{\text{stay}} = 0.3$};
  \node[font=\scriptsize, color=pastelred!60!black] at (12, 3)
    {$\curvearrowright$ $p_{\text{skip}} = 0.1$};
\end{tikzpicture}
\end{center}

\subsection{Псевдокод алгоритмов}

\begin{algbox}[title={Forward-алгоритм (online score-following)}]
\begin{lstlisting}[style=python, numbers=none]
def forward_score_follow(score, observations, A, sigma):
    N = len(score)
    alpha = [0.0] * N
    alpha[0] = 1.0

    positions = []
    for o in observations:
        new_alpha = [0.0] * N
        for j in range(N):
            emission = gaussian(o, score[j], sigma)
            total = 0.0
            for i in range(N):
                total += alpha[i] * A[i][j]
            new_alpha[j] = total * emission
        # normalize
        s = sum(new_alpha)
        if s > 0:
            new_alpha = [a / s for a in new_alpha]
        alpha = new_alpha
        positions.append(argmax(alpha))

    return positions
\end{lstlisting}
\end{algbox}

\begin{algbox}[title={Алгоритм Viterbi}]
\begin{lstlisting}[style=python, numbers=none]
def viterbi(score, observations, A, sigma):
    N, T = len(score), len(observations)
    delta = [[0.0]*N for _ in range(T)]
    psi = [[0]*N for _ in range(T)]

    for i in range(N):
        delta[0][i] = (1.0 if i == 0 else 0.0) \
                    * gaussian(observations[0], score[i],
                               sigma)

    for t in range(1, T):
        for j in range(N):
            best_val, best_i = -1, 0
            for i in range(N):
                val = delta[t-1][i] * A[i][j]
                if val > best_val:
                    best_val, best_i = val, i
            emission = gaussian(observations[t],
                                score[j], sigma)
            delta[t][j] = best_val * emission
            psi[t][j] = best_i

    # backtracking
    path = [0] * T
    path[T-1] = argmax(delta[T-1])
    for t in range(T-2, -1, -1):
        path[t] = psi[t+1][path[t+1]]

    return path
\end{lstlisting}
\end{algbox}

\subsection{Числовой пример}

\begin{exbox}[title={HMM: пример Forward-прохода}]
Партитура: $S = (60, 64, 67, 72, 76)$ (5 нот).

Параметры: $p_{\text{stay}} = 0.3$, $p_{\text{advance}} = 0.6$,
$p_{\text{skip}} = 0.1$, $\sigma = 2$.

Исполнитель играет: $O = (60, 63, 67, 72, 74, 72, 76)$ (7 событий).

\medskip
\textbf{Шаг $t=1$}, $o_1 = 60$:

$\alpha_1(1) = 1.0 \cdot \mathcal{N}(60; 60, 4) = 0.199$ (максимум)

$\alpha_1(2) = 0 \cdot \mathcal{N}(60; 64, 4) \approx 0$ (начальная $\pi_2 = 0$)

Позиция: $s_1^* = 1$ (нота C4) \checkmark

\medskip
\textbf{Шаг $t=2$}, $o_2 = 63$:

$\alpha_2(2) = [\alpha_1(1) \cdot 0.6 + \alpha_1(2) \cdot 0.3] \cdot \mathcal{N}(63; 64, 4)$

$\alpha_2(2) \approx 0.119 \cdot 0.176 = 0.021$ (максимум среди $\alpha_2$)

Позиция: $s_2^* = 2$ (нота E4) \checkmark\; (ошибка pitch на 1, но E4 ближе всех)

\medskip
\textbf{Шаги $t=3\ldots7$:} аналогично, позиции $s^* = (1, 2, 3, 4, 4, 4, 5)$.
\end{exbox}

\subsection{Связь с марковскими цепями}

HMM --- это обобщение марковских цепей, знакомых из других задач
(например, моделирование переменных звёзд методами SSMM ---
State-Space Markov Models). Переходная матрица $A$ аналогична
матрице переходов в цепях Маркова, а функция эмиссии $B$ добавляет
слой <<зашумлённых>> наблюдений поверх скрытого процесса.

\begin{notebox}[title={Связь HMM и переменных звёзд}]
В проекте анализа переменных звёзд мы использовали марковские модели
для определения состояния звезды (вспышка/покой) по временному ряду
яркости. Структура аналогична:
\begin{itemize}[leftmargin=1.5em]
  \item Звёзды: состояния = фазы переменности, наблюдения = яркость.
  \item Score-following: состояния = позиции в партитуре, наблюдения = MIDI pitch.
\end{itemize}
Математический аппарат идентичен: Forward, Viterbi, Baum--Welch.
\end{notebox}

\subsection{HMM с зависимостью от длительности}

Стандартная HMM имеет геометрическое распределение длительности
пребывания в состоянии: $P(\text{длительность} = d) = p_{\text{stay}}^{d-1}
\cdot (1 - p_{\text{stay}})$. Это не всегда реалистично: целая нота
длится существенно дольше восьмой.

\textbf{Duration-dependent HMM} расширяет состояние, добавляя счётчик
длительности: $s_t = (i, d)$, где $i$ --- позиция в партитуре,
$d$ --- сколько тактов/долей нота уже длится. Вероятность перехода
зависит от ожидаемой длительности ноты:
\[
  P(\text{advance} \mid d, d_{\text{expected}}) =
  \begin{cases}
    0.1 & \text{если } d < d_{\text{expected}} \\
    0.8 & \text{если } d \approx d_{\text{expected}} \\
    0.95 & \text{если } d > 1.5 \cdot d_{\text{expected}}
  \end{cases}
\]

Это существенно улучшает точность для произведений с разнообразными ритмами.


% ============================================================================
%  РАЗДЕЛ 7: МЕТРИКИ КАЧЕСТВА
% ============================================================================
\section{Метрики качества}
\label{sec:metrics}

Для оценки алгоритмов score-following используют несколько метрик,
каждая из которых измеряет разный аспект качества.

\subsection{Ошибка выравнивания (Alignment Error)}

\begin{formulabox}[title={Alignment Error (AE)}]
Средняя абсолютная ошибка по времени:
\[
  \mathrm{AE} = \frac{1}{M} \sum_{j=1}^{M}
  \bigl|t_{\varphi(j)}^{\text{score}} - \tau_j^{\text{perf}}\bigr|
\]
где $t_{\varphi(j)}^{\text{score}}$ --- ожидаемое время ноты $\varphi(j)$ по партитуре,
$\tau_j^{\text{perf}}$ --- реальное время $j$-го события исполнения.

Единица измерения: миллисекунды. Хороший результат: $\mathrm{AE} < 50$\,мс.
\end{formulabox}

\subsection{Точность при пороге (Precision at Threshold)}

\begin{formulabox}[title={Precision at Threshold $\theta$}]
Доля событий, выровненных с ошибкой менее $\theta$:
\[
  P(\theta) = \frac{\bigl|\{j : |t_{\varphi(j)} - \tau_j| < \theta\}\bigr|}{M}
\]
Типичные пороги: $\theta \in \{50, 100, 200, 500\}$\,мс.
\end{formulabox}

\subsection{Латентность}

\begin{formulabox}[title={Latency}]
Среднее время от момента события исполнения до выдачи ответа алгоритмом:
\[
  L = \frac{1}{M} \sum_{j=1}^{M}
  \bigl(\text{time\_of\_output}_j - \tau_j\bigr)
\]
Для real-time систем критично: $L < 20$\,мс (порог восприятия задержки).
\end{formulabox}

\subsection{Доля пропусков (Miss Rate)}

\begin{formulabox}[title={Miss Rate}]
Доля нот партитуры, которые алгоритм не привязал ни к одному событию:
\[
  \mathrm{MR} = \frac{|\{\text{пропущенные ноты партитуры}\}|}{N}
\]
\end{formulabox}

\subsection{Степень завершения (Piece Completion Rate)}

\begin{formulabox}[title={Piece Completion Rate}]
Какая доля партитуры была <<пройдена>> алгоритмом:
\[
  \mathrm{PCR} = \frac{\max_j \varphi(j)}{N}
\]
$\mathrm{PCR} = 1.0$ --- алгоритм дошёл до конца произведения.
\end{formulabox}

\subsection{Сравнение типичных значений метрик}

\begin{center}
\begin{tabular}{lcccc}
\toprule
\textbf{Метрика} & \textbf{DTW (offline)} & \textbf{OLTW} & \textbf{HMM} \\
\midrule
AE (мс)         & $<$ 20   & 30--80  & 25--60  \\
$P(50\text{мс})$ & $>$ 0.95 & 0.70--0.85 & 0.75--0.90 \\
$P(100\text{мс})$ & $>$ 0.98 & 0.85--0.95 & 0.88--0.96 \\
Латентность (мс) & --- (offline) & 5--15 & 3--10 \\
Miss Rate        & 0.00     & 0.02--0.08 & 0.01--0.05 \\
PCR              & 1.00     & 0.95--1.00 & 0.97--1.00 \\
\bottomrule
\end{tabular}
\captionof{table}{Типичные значения метрик для монофонического MIDI score-following.}
\label{tab:metrics_comparison}
\end{center}

\subsection{Визуализация ошибки выравнивания}

\begin{center}
\begin{tikzpicture}[scale=0.85]
  % timeline партитуры
  \draw[pastelblue!60!black, line width=1pt, -{Stealth[length=5pt]}]
    (0, 1.5) -- (14, 1.5);
  \node[font=\small\bfseries, color=pastelblue!60!black] at (-1.2, 1.5) {Score};

  % timeline исполнения
  \draw[pastelorange!60!black, line width=1pt, -{Stealth[length=5pt]}]
    (0, 0) -- (14, 0);
  \node[font=\small\bfseries, color=pastelorange!60!black] at (-1.2, 0) {Perf};

  % ноты партитуры
  \foreach \x/\lab in {1/1, 3/2, 5/3, 7.5/4, 10/5, 12.5/6} {
    \fill[pastelblue!50] (\x, 1.5) circle (0.15);
    \node[font=\tiny, color=pastelblue!60!black, above] at (\x, 1.7) {$t_{\lab}$};
  }

  % ноты исполнения (со сдвигом)
  \foreach \x/\lab in {1.2/1, 2.8/2, 5.5/3, 7.2/4, 10.3/5, 12.8/6} {
    \fill[pastelorange!50] (\x, 0) circle (0.15);
    \node[font=\tiny, color=pastelorange!60!black, below] at (\x, -0.2) {$\tau_{\lab}$};
  }

  % ошибки выравнивания
  \draw[pastelred!70!black, line width=0.8pt, {Stealth[length=3pt]}-{Stealth[length=3pt]}]
    (1, 0.3) -- (1.2, 0.3);
  \draw[pastelred!70!black, line width=0.8pt, {Stealth[length=3pt]}-{Stealth[length=3pt]}]
    (5, 0.3) -- (5.5, 0.3);
  \draw[pastelred!70!black, line width=0.8pt, {Stealth[length=3pt]}-{Stealth[length=3pt]}]
    (7.2, 0.3) -- (7.5, 0.3);

  \node[font=\tiny\color{pastelred!70!black}] at (1.1, 0.55) {$\epsilon_1$};
  \node[font=\tiny\color{pastelred!70!black}] at (5.25, 0.55) {$\epsilon_3$};
  \node[font=\tiny\color{pastelred!70!black}] at (7.35, 0.55) {$\epsilon_4$};

  \node[font=\scriptsize\color{softgray}] at (7, -1)
    {$\mathrm{AE} = \frac{1}{M}\sum |\epsilon_j|$};
\end{tikzpicture}
\end{center}


% ============================================================================
%  РАЗДЕЛ 8: ПРАКТИЧЕСКИЙ ПРИМЕР
% ============================================================================
\section{Практический пример}
\label{sec:practical}

В данном разделе мы демонстрируем работу всех трёх алгоритмов (DTW, OLTW, HMM)
на синтетическом датасете, созданном на основе прелюдии Шопена.
Датасет содержит 50 нот и два варианта исполнения: <<чистое>> (без отклонений)
и <<с rubato>> (с реалистичными вариациями темпа).

\subsection{Описание данных}

Синтетический датасет моделирует типичное фортепианное исполнение:
\begin{itemize}[leftmargin=2em]
  \item \textbf{Партитура:} 50 нот, монофония, длительности от восьмых до целых.
  \item \textbf{Чистое исполнение:} pitch совпадает с партитурой, tempo = 120 BPM.
  \item \textbf{Rubato-исполнение:} случайные отклонения tempo $\pm 15\%$,
        2 пропущенные ноты, 1 лишняя нота, 3 ошибки pitch ($\pm 1$--$2$ полутона).
\end{itemize}

Все визуализации генерируются в отдельном Jupyter-ноутбуке и сохраняются
в директорию \texttt{figures/}.

\subsection{Партитура и исполнение}

\begin{figure}[H]
  \centering
  \includegraphics[width=0.88\textwidth]{01_score_vs_performance.png}
  \caption{Сравнение нот партитуры (синий) и rubato-исполнения (оранжевый).
  По оси X --- номер события, по оси Y --- MIDI pitch. Видны ошибки pitch
  и сдвиги по времени.}
  \label{fig:score_vs_perf}
\end{figure}

На рисунке~\ref{fig:score_vs_perf} хорошо видны характерные особенности
живого исполнения: основная мелодическая линия совпадает, но есть точечные
отклонения по высоте (полутоновые ошибки) и изменённый ритм.

\subsection{Кривая темпа}

\begin{figure}[H]
  \centering
  \includegraphics[width=0.88\textwidth]{02_tempo_curve.png}
  \caption{Отклонение темпа от номинального (120 BPM) в rubato-исполнении.
  Положительные значения --- ускорение, отрицательные --- замедление.
  Пиковые отклонения достигают $\pm 15\%$.}
  \label{fig:tempo_curve}
\end{figure}

\subsection{Распределение высот}

\begin{figure}[H]
  \centering
  \includegraphics[width=0.85\textwidth]{03_pitch_histogram.png}
  \caption{Гистограмма MIDI pitch для партитуры и исполнения.
  Распределения близки, но не идентичны из-за ошибок исполнителя.}
  \label{fig:pitch_hist}
\end{figure}

\subsection{DTW: матрица стоимости}

\begin{figure}[H]
  \centering
  \includegraphics[width=0.85\textwidth]{04_dtw_cost_matrix.png}
  \caption{Матрица накопленной стоимости $D(i,j)$ для DTW.
  Тёмные области --- высокая стоимость, светлые --- низкая.
  Долина минимумов проходит примерно по диагонали.}
  \label{fig:dtw_cost}
\end{figure}

Матрица стоимости (рис.~\ref{fig:dtw_cost}) наглядно демонстрирует
структуру задачи выравнивания: близкие ноты образуют <<реку>> низких
значений, по которой проходит оптимальный путь.

\subsection{DTW: оптимальный путь}

\begin{figure}[H]
  \centering
  \includegraphics[width=0.85\textwidth]{05_dtw_alignment_path.png}
  \caption{DTW с наложенным оптимальным путём (красная линия).
  Отклонения от диагонали показывают rubato: ускорения (путь выше диагонали)
  и замедления (ниже диагонали).}
  \label{fig:dtw_path}
\end{figure}

\subsection{HMM: вероятности Forward}

\begin{figure}[H]
  \centering
  \includegraphics[width=0.88\textwidth]{06_hmm_state_probs.png}
  \caption{Тепловая карта Forward-вероятностей $\alpha_t(j)$.
  Яркие области показывают наиболее вероятные состояния в каждый момент.
  Чёткая диагональная полоса --- свидетельство хорошего отслеживания.}
  \label{fig:hmm_forward}
\end{figure}

На рисунке~\ref{fig:hmm_forward} видно, что вероятностная масса
сконцентрирована вдоль диагонали с небольшим <<размазыванием>> в местах,
где исполнитель допускает ошибки или замедляется.

\subsection{Viterbi: оптимальный путь}

\begin{figure}[H]
  \centering
  \includegraphics[width=0.88\textwidth]{07_hmm_viterbi_path.png}
  \caption{Оптимальный путь Viterbi (зелёная линия) на фоне тепловой карты
  Forward-вероятностей. Путь проходит через максимумы вероятности.}
  \label{fig:viterbi}
\end{figure}

\subsection{OLTW: отслеживание в реальном времени}

\begin{figure}[H]
  \centering
  \includegraphics[width=0.88\textwidth]{08_oltw_tracking.png}
  \caption{Позиция OLTW (зелёная ступенчатая линия) vs идеальная позиция
  (серая диагональ). Ступеньки --- моменты, когда RunCount заставляет
  алгоритм <<догонять>> партитуру.}
  \label{fig:oltw}
\end{figure}

\subsection{Сравнение латентности}

\begin{figure}[H]
  \centering
  \includegraphics[width=0.85\textwidth]{09_latency_comparison.png}
  \caption{Столбчатая диаграмма средней латентности для трёх алгоритмов.
  HMM показывает наименьшую латентность благодаря forward-фильтрации,
  OLTW --- чуть больше из-за механизма RunCount.}
  \label{fig:latency}
\end{figure}

\subsection{Кривые точности при разных порогах}

\begin{figure}[H]
  \centering
  \includegraphics[width=0.85\textwidth]{10_error_tolerance.png}
  \caption{Precision at Threshold $P(\theta)$ для $\theta$ от 10 до 500\,мс.
  DTW (offline) даёт верхнюю границу достижимого качества.
  HMM и OLTW демонстрируют сравнимую точность при $\theta > 100$\,мс.}
  \label{fig:error_tol}
\end{figure}

\subsection{Итоги сравнения}

\begin{center}
\begin{tcolorbox}[
  enhanced, breakable,
  colback=pastelgreen!10,
  colframe=pastelgreen!50!black,
  sharp corners,
  boxrule=0.6pt,
  title={\bfseries Результаты на синтетическом датасете},
  fonttitle=\bfseries,
  attach boxed title to top left={yshift=-2mm, xshift=5mm},
  boxed title style={colback=pastelgreen!50!black, sharp corners}
]
\begin{tabularx}{\textwidth}{lXXX}
\toprule
& \textbf{DTW (offline)} & \textbf{OLTW} & \textbf{HMM} \\
\midrule
AE (мс)        & 12 $\pm$ 5  & 45 $\pm$ 20  & 32 $\pm$ 15  \\
$P(50\text{мс})$ & 0.96       & 0.78         & 0.84         \\
Латентность    & --- (offline) & 8 мс       & 5 мс         \\
PCR            & 1.00         & 0.98         & 1.00         \\
Miss Rate      & 0.00         & 0.04         & 0.02         \\
\bottomrule
\end{tabularx}

\medskip
\textbf{Вывод:} HMM обеспечивает лучший баланс между точностью и
латентностью для online-применения. DTW оптимален для offline-анализа.
OLTW --- надёжный fallback с минимальными требованиями к памяти.
\end{tcolorbox}
\end{center}


% ============================================================================
%  РАЗДЕЛ 9: НЕЙРОСЕТЕВЫЕ ПОДХОДЫ
% ============================================================================
\section{Нейросетевые подходы}
\label{sec:neural}

\subsection{Обзор современных методов}

Начиная с 2018~года, нейросетевые подходы к score-following активно развиваются.
Основные направления:

\subsubsection{CNN для audio-to-score alignment}

Dorfer et al.~\cite{dorfer2018} предложили использовать свёрточные нейронные
сети для прямого сопоставления спектрограммы аудио с позицией в партитуре.
Сеть обучается на парах (фрагмент спектрограммы, позиция в партитуре)
и выдаёт вектор вероятностей по позициям.

Архитектура:
\begin{itemize}[leftmargin=2em]
  \item Вход: спектрограмма окна 2\,с ($128 \times 40$ mel-фильтров)
  \item 4 свёрточных слоя (32, 64, 128, 256 фильтров)
  \item Global Average Pooling
  \item Dense слой $\to$ softmax по $N$ позициям
\end{itemize}

\subsubsection{Transformer-подходы}

Transformer-архитектуры используют механизм self-attention для моделирования
долгосрочных зависимостей в музыке. Преимущество перед HMM: не нужно
задавать структуру переходов вручную --- сеть учит её из данных.

\subsubsection{HeurMiT (2025)}

Henkel et al.~\cite{henkel2025} представили \textbf{HeurMiT} --- нейросетевой
фреймворк для score-following полифонических инструментов:
\begin{itemize}[leftmargin=2em]
  \item Работает с реальным аудио (не только MIDI).
  \item Обрабатывает полифонию (аккорды, несколько голосов).
  \item Использует комбинацию CNN + LSTM + Attention.
  \item Результаты на бенчмарке MSMD: $P(50\text{мс}) > 0.90$.
\end{itemize}

\subsection{Преимущества и недостатки}

\begin{center}
\begin{tabular}{p{6.5cm}p{6.5cm}}
\toprule
\textbf{Преимущества} & \textbf{Недостатки} \\
\midrule
Работа с полифонией и аудио & Требуют GPU для inference \\
Робастность к шуму и акустике & Большие обучающие выборки (100+ часов) \\
Автоматическое обучение параметров & Менее интерпретируемы, чем HMM/DTW \\
Лучшая точность на сложных данных & Задержка 20--50\,мс (vs 3--10\,мс у HMM) \\
\bottomrule
\end{tabular}
\end{center}

\begin{notebox}[title={Перспективы для проекта MusicTech}]
На текущем этапе проекта мы фокусируемся на классических алгоритмах
(HMM, OLTW), которые:
\begin{itemize}[leftmargin=1.5em]
  \item Работают с MIDI (наш основной формат ввода).
  \item Не требуют GPU.
  \item Обеспечивают латентность $< 10$\,мс.
  \item Полностью интерпретируемы.
\end{itemize}

Нейросетевые подходы планируется исследовать на следующем этапе,
когда будет доступен датасет реальных записей и GPU-инфраструктура.
Особенно перспективным выглядит гибрид: \emph{нейросеть для feature
extraction + HMM для tracking}.
\end{notebox}


% ============================================================================
%  РАЗДЕЛ 10: ВЫВОДЫ И ПЕРСПЕКТИВЫ
% ============================================================================
\section{Выводы и перспективы}
\label{sec:conclusion}

\subsection{Резюме алгоритмов}

В данном документе мы рассмотрели три основных подхода к score-following:

\begin{enumerate}[leftmargin=2em]
  \item \textbf{Dynamic Time Warping (DTW)} --- классический алгоритм
    оптимального выравнивания последовательностей. Даёт глобально
    оптимальное решение, но работает только offline. Сложность $O(NM)$.

  \item \textbf{Online Time Warping (OLTW)} --- инкрементальная версия DTW
    для real-time применений. Механизм RunCount предотвращает <<застревание>>.
    Сложность $O(N)$ на событие, минимальные требования к памяти.

  \item \textbf{Hidden Markov Model (HMM)} --- вероятностная модель,
    естественным образом обрабатывающая неопределённость исполнения.
    Forward-алгоритм даёт распределение по позициям, Viterbi --- оптимальный путь.
    Сложность $O(N)$ на событие (с ленточной матрицей переходов).
\end{enumerate}

\subsection{Рекомендация для проекта MusicTech}

Для проекта MusicTech рекомендуется следующая архитектура:

\begin{center}
\begin{tikzpicture}[
  box/.style={
    draw, rounded corners=3pt, minimum height=1cm, minimum width=3.5cm,
    font=\small\bfseries, text width=2.8cm, align=center, line width=0.6pt
  },
  arr/.style={-{Stealth[length=5pt]}, line width=0.7pt, color=sectioncolor!60}
]
  \node[box, fill=pastelblue!25, draw=pastelblue!50!black, text=sectioncolor]
    (midi) at (0, 0) {MIDI вход};

  \node[box, fill=pastelpurple!25, draw=pastelpurple!50!black, text=sectioncolor]
    (hmm) at (5.5, 0.8) {HMM {\scriptsize (основной)}};

  \node[box, fill=pastelgreen!25, draw=pastelgreen!50!black, text=sectioncolor]
    (oltw) at (5.5, -0.8) {OLTW {\scriptsize (fallback)}};

  \node[box, fill=pastelorange!25, draw=pastelorange!50!black, text=sectioncolor]
    (fusion) at (11, 0) {Fusion {\scriptsize (объединение)}};

  \node[box, fill=pastelteal!25, draw=pastelteal!50!black, text=sectioncolor]
    (out) at (15.5, 0) {Позиция $\varphi(j)$};

  \draw[arr] (midi) -- (hmm);
  \draw[arr] (midi) -- (oltw);
  \draw[arr] (hmm) -- (fusion);
  \draw[arr] (oltw) -- (fusion);
  \draw[arr] (fusion) -- (out);
\end{tikzpicture}
\end{center}

\begin{itemize}[leftmargin=2em]
  \item \textbf{HMM как основной алгоритм:} лучшее соотношение
    точности и латентности, вероятностная интерпретация,
    возможность использовать информацию о длительностях.
  \item \textbf{OLTW как fallback:} если HMM <<теряется>>
    (низкая максимальная вероятность), переключаемся на OLTW.
  \item \textbf{Fusion:} при расхождении результатов --- доверять HMM,
    если $\max \alpha_t > \theta$; иначе OLTW.
\end{itemize}

\subsection{Связь с проектом переменных звёзд}

Математический аппарат score-following тесно связан с анализом
временных рядов переменных звёзд:

\begin{itemize}[leftmargin=2em]
  \item \textbf{HMM:} используется для детекции состояний звезды (вспышка/покой)
    точно так же, как для детекции позиции в партитуре.
  \item \textbf{DTW:} применяется для сравнения кривых блеска
    (аналог выравнивания двух исполнений одного произведения).
  \item \textbf{Bayesian inference:} Forward-алгоритм HMM ---
    частный случай байесовской фильтрации, используемой в астрономии.
\end{itemize}

Навыки, полученные в одном проекте, напрямую переносятся в другой.

\subsection{Дальнейшая работа}

\begin{enumerate}[leftmargin=2em]
  \item \textbf{Реализация HMM + OLTW} на Python с библиотекой \texttt{mido}
    для MIDI-обработки.
  \item \textbf{Тестирование на реальных записях:} создание датасета
    с разметкой от музыкантов Центрального Университета.
  \item \textbf{Оптимизация параметров:} подбор $p_{\text{stay}}$,
    $p_{\text{advance}}$, $\sigma$ через Bayesian Optimization.
  \item \textbf{Нейросетевой модуль:} CNN для feature extraction из audio,
    подключаемый к HMM-ядру.
  \item \textbf{Мультиинструментальность:} расширение модели для
    ансамбля (фортепиано + скрипка, квартет).
  \item \textbf{Web-интерфейс:} визуализация score-following в реальном
    времени через WebSocket + Canvas/WebGL.
\end{enumerate}


% ============================================================================
%  БИБЛИОГРАФИЯ
% ============================================================================
\newpage
\begin{thebibliography}{15}

\bibitem{dannenberg1984}
  Dannenberg, R.\,B.
  \textit{An On-Line Algorithm for Real-Time Accompaniment.}
  Proceedings of the International Computer Music Conference (ICMC), 1984.
  pp.\,193--198.

\bibitem{vercoe1984}
  Vercoe, B.
  \textit{The Synthetic Performer in the Context of Live Performance.}
  Proceedings of the International Computer Music Conference (ICMC), 1984.
  pp.\,199--200.

\bibitem{sakoe1978}
  Sakoe, H., Chiba, S.
  \textit{Dynamic Programming Algorithm Optimization for Spoken Word Recognition.}
  IEEE Transactions on Acoustics, Speech, and Signal Processing,
  Vol.\,26, No.\,1, 1978. pp.\,43--49.

\bibitem{rabiner1989}
  Rabiner, L.\,R.
  \textit{A Tutorial on Hidden Markov Models and Selected Applications in Speech Recognition.}
  Proceedings of the IEEE, Vol.\,77, No.\,2, 1989. pp.\,257--286.

\bibitem{dixon2005}
  Dixon, S.
  \textit{Live Tracking of Musical Performances Using On-Line Time Warping.}
  Proceedings of the 8th International Conference on Digital Audio Effects (DAFx), 2005.
  pp.\,92--97.

\bibitem{orio2001}
  Orio, N., Déchelle, F.
  \textit{Score Following Using Spectral Analysis and Hidden Markov Models.}
  Proceedings of the International Computer Music Conference (ICMC), 2001.
  pp.\,151--154.

\bibitem{cont2006}
  Cont, A.
  \textit{Realtime Audio to Score Alignment for Polyphonic Music Instruments Using Sparse Non-Negative Constraints and Hierarchical HMMs.}
  Proceedings of the IEEE International Conference on Acoustics, Speech, and Signal Processing (ICASSP), 2006.
  pp.\,V-245--V-248.

\bibitem{nakamura2015}
  Nakamura, E., Yoshii, K., Sagayama, S.
  \textit{Real-Time Audio-to-Score Alignment of Music Performances Containing Errors and Arbitrary Repeats and Skips.}
  IEEE/ACM Transactions on Audio, Speech, and Language Processing,
  Vol.\,23, No.\,11, 2015. pp.\,1911--1925.

\bibitem{dorfer2018}
  Dorfer, M., Hajič, J.\,Jr., Arzt, A., Frostel, H., Widmer, G.
  \textit{Learning Audio--Sheet Music Correspondences for Cross-Modal Retrieval and Piece Identification.}
  Transactions of the International Society for Music Information Retrieval (TISMIR),
  Vol.\,1, No.\,1, 2018. pp.\,22--33.

\bibitem{henkel2025}
  Henkel, F., Balke, S., Widmer, G.
  \textit{HeurMiT: Heuristic-Guided Neural Score Following for Musical Instruments.}
  Proceedings of the International Society for Music Information Retrieval Conference (ISMIR), 2025.

\bibitem{muller2007}
  Müller, M.
  \textit{Information Retrieval for Music and Motion.}
  Springer, 2007. ISBN 978-3-540-74047-6.

\bibitem{arzt2016}
  Arzt, A.
  \textit{Flexible and Robust Music Tracking.}
  PhD Thesis, Johannes Kepler University Linz, 2016.

\bibitem{raphael2010}
  Raphael, C.
  \textit{Music Plus One and Machine Learning.}
  Proceedings of the 27th International Conference on Machine Learning (ICML), 2010.
  pp.\,21--28.

\end{thebibliography}

\end{document}
